\documentclass[11pt,a4paper]{article}
\usepackage[utf8]{inputenc}
\usepackage[T1]{fontenc}
\usepackage[margin=1in]{geometry}
\usepackage{xcolor}
\usepackage{hyperref}
\definecolor{primarycolor}{RGB}{33,37,41}
\definecolor{secondarycolor}{RGB}{0,102,204}
\hypersetup{colorlinks=true,linkcolor=secondarycolor,urlcolor=secondarycolor}
\pagestyle{empty}
\begin{document}
\noindent Dear Tessel Cargo Software Hiring Manager,

I am applying for the Backend Engineer (Python) role. I have built backend services professionally in Java and Kotlin since 2019, and I am keen to bring that experience to a Python-focused team. My Python experience so far has been scripting, not production services, but I am interested in making the move to a Python stack.

At Tejo Insurance Tech, I designed a versioned REST API used by 12 broker partners for around 400,000 quote requests per month. I have also worked directly on the reliability of event-driven services: after implementing a transactional outbox and idempotency keys for Kafka consumers, duplicate claims fell from around 40 per week to zero over the following quarter. These projects have given me relevant experience with REST APIs, PostgreSQL, Kafka, automated testing and the care required when processing events reliably.

I also participate in a monthly on-call rota and handled a major Kafka lag incident in March 2024; the incident write-up was praised. I would bring that operational experience, along with hands-on use of Kubernetes on EKS, Docker, AWS S3 and SQS, to a team responsible for running its own tracking services. Tessel’s focus on dependable live shipment tracking and event streams makes this role a particularly relevant next step for my backend experience.

Thank you for considering my application. I would welcome the opportunity to discuss how my experience could contribute to the Routing & Tracking team.

Best regards,\\
Daniel Okafor-Reyes
\end{document}